\capituloPregunta{¿Cómo se convierten las presentaciones en prácticas y capítulos?}

        Una presentacion aislada puede servir para exponer un tema, pero no siempre basta para que un estudiante aprenda de manera autonoma. En este manual, cada presentacion compatible con ADE debe dejar tres rastros verificables: una practica, una seccion narrativa del manual y una version enriquecida para estudiar.

        \begin{idea}
        La regla operativa es simple: si existe una presentacion y falta la practica, se crea la practica; si existe la practica y la presentacion contiene mejores ejemplos, la practica y el manual se enriquecen; si falta la narrativa, se escribe antes de publicar.
        \end{idea}

        \section{Presentaciones integradas por el sistema}

        \begin{itemize}
        \item \textbf{Muestreo e indice de Gini}: practica generada, seccion de manual y presentacion enriquecida.
\item \textbf{De datos continuos a diseno Taguchi con Python}: practica generada, seccion de manual y presentacion enriquecida.
\item \textbf{MANOVA con respuestas fisiologicas en plantas}: practica generada, seccion de manual y presentacion enriquecida.
\item \textbf{Rendimiento termico en procesadores}: practica generada, seccion de manual y presentacion enriquecida.
\item \textbf{ANCOVA con tres variables explicativas}: practica generada, seccion de manual y presentacion enriquecida.
        \end{itemize}

        \section{Presentacion convertida en practica: Muestreo e indice de Gini}

Esta seccion nace de una presentacion existente y se incorpora al manual para que el material no permanezca aislado. La presentacion deja de ser un apoyo oral y se transforma en una ruta de estudio: problema, variables, datos, analisis, decision y actividad.

\begin{idea}
Comprender como distintos metodos de muestreo modifican la representacion de una realidad social.
\end{idea}

\subsection{Evidencia didactica extraida}

\begin{itemize}
\item Por ello, la estadística utiliza técnicas de muestreo que permiten
\end{itemize}

\subsection{Como entra al manual}

El tema se integra al capitulo relacionado con \textbf{muestreo}. Si el estudiante solo observa la presentacion, ve una secuencia visual. En el manual, esa misma secuencia se expande con la pregunta de ingenieria, la razon del metodo, la interpretacion y los limites de la conclusion.

\subsection{Practica asociada}

La practica generada pide al estudiante transformar la presentacion en una decision: definir variables, organizar datos, graficar, aplicar el metodo y redactar una conclusion prudente. Si no existia practica, el sistema crea una primera version editable. Si existia una practica relacionada, esta se usa para enriquecer la narrativa y no como bloque aislado.

\subsection{Actividad de evaluacion}

Explique que parte de la presentacion corresponde a contexto, cual corresponde a diseno, cual corresponde a evidencia y cual corresponde a decision. Si falta alguno de estos elementos, proponga como completarlo.

\section{Presentacion convertida en practica: De datos continuos a diseno Taguchi con Python}

Esta seccion nace de una presentacion existente y se incorpora al manual para que el material no permanezca aislado. La presentacion deja de ser un apoyo oral y se transforma en una ruta de estudio: problema, variables, datos, analisis, decision y actividad.

\begin{idea}
Convertir factores continuos en niveles discretos para estudiar robustez y reproducibilidad.
\end{idea}

\subsection{Evidencia didactica extraida}

\begin{itemize}
\item Una metodología práctica para convertir datos experimentales continuos en diseños
\item robustos tipo Taguchi, con aplicaciones en ingeniería y estadística experimental.
\item Explicar cómo transformar variables continuas en
\item niveles discretos Taguchi mediante técnicas de
\end{itemize}

\subsection{Como entra al manual}

El tema se integra al capitulo relacionado con \textbf{taguchi}. Si el estudiante solo observa la presentacion, ve una secuencia visual. En el manual, esa misma secuencia se expande con la pregunta de ingenieria, la razon del metodo, la interpretacion y los limites de la conclusion.

\subsection{Practica asociada}

La practica generada pide al estudiante transformar la presentacion en una decision: definir variables, organizar datos, graficar, aplicar el metodo y redactar una conclusion prudente. Si no existia practica, el sistema crea una primera version editable. Si existia una practica relacionada, esta se usa para enriquecer la narrativa y no como bloque aislado.

\subsection{Actividad de evaluacion}

Explique que parte de la presentacion corresponde a contexto, cual corresponde a diseno, cual corresponde a evidencia y cual corresponde a decision. Si falta alguno de estos elementos, proponga como completarlo.

\section{Presentacion convertida en practica: MANOVA con respuestas fisiologicas en plantas}

Esta seccion nace de una presentacion existente y se incorpora al manual para que el material no permanezca aislado. La presentacion deja de ser un apoyo oral y se transforma en una ruta de estudio: problema, variables, datos, analisis, decision y actividad.

\begin{idea}
Decidir cuando varias respuestas correlacionadas deben analizarse como sistema.
\end{idea}

\subsection{Evidencia didactica extraida}

\begin{itemize}
\item Análisis Multivariado de la Varianza con Tres Variables Dependientes
\item Evaluar si los factores experimentales (luz, riego y fertilizante) influyen de manera
\item significativa en tres variables de respuesta fisiológica de las plantas:
\item Tipo: Factorial 3 × 3 × 3 reducido a 10 tratamientos
\end{itemize}

\subsection{Como entra al manual}

El tema se integra al capitulo relacionado con \textbf{manova}. Si el estudiante solo observa la presentacion, ve una secuencia visual. En el manual, esa misma secuencia se expande con la pregunta de ingenieria, la razon del metodo, la interpretacion y los limites de la conclusion.

\subsection{Practica asociada}

La practica generada pide al estudiante transformar la presentacion en una decision: definir variables, organizar datos, graficar, aplicar el metodo y redactar una conclusion prudente. Si no existia practica, el sistema crea una primera version editable. Si existia una practica relacionada, esta se usa para enriquecer la narrativa y no como bloque aislado.

\subsection{Actividad de evaluacion}

Explique que parte de la presentacion corresponde a contexto, cual corresponde a diseno, cual corresponde a evidencia y cual corresponde a decision. Si falta alguno de estos elementos, proponga como completarlo.

\section{Presentacion convertida en practica: Rendimiento termico en procesadores}

Esta seccion nace de una presentacion existente y se incorpora al manual para que el material no permanezca aislado. La presentacion deja de ser un apoyo oral y se transforma en una ruta de estudio: problema, variables, datos, analisis, decision y actividad.

\begin{idea}
Equilibrar temperatura y rendimiento en computacion intensiva mediante diseno factorial y respuestas multiples.
\end{idea}

\subsection{Evidencia didactica extraida}

\begin{itemize}
\item Aplicación de un Diseño Factorial 3\textasciicircum{}3 y MANOVA para la Optimización de Energía y Ventilación.
\item El ANOVA para la variable Puntos muestra que los Hilos tienen un valor F inmenso (15,732), siendo el factor dominante.
\item Análisis Térmico: ANOVA y Efectividad del Ventilador.
\item Para la Temperatura , el ANOVA revela algo distinto: el factor Energía NO es significativo (p=0.20, vean la tabla), lo que significa que el CPU se calienta igual en modo 'Ahorro' o 'Alto'.
\end{itemize}

\subsection{Como entra al manual}

El tema se integra al capitulo relacionado con \textbf{manova factorial}. Si el estudiante solo observa la presentacion, ve una secuencia visual. En el manual, esa misma secuencia se expande con la pregunta de ingenieria, la razon del metodo, la interpretacion y los limites de la conclusion.

\subsection{Practica asociada}

La practica generada pide al estudiante transformar la presentacion en una decision: definir variables, organizar datos, graficar, aplicar el metodo y redactar una conclusion prudente. Si no existia practica, el sistema crea una primera version editable. Si existia una practica relacionada, esta se usa para enriquecer la narrativa y no como bloque aislado.

\subsection{Actividad de evaluacion}

Explique que parte de la presentacion corresponde a contexto, cual corresponde a diseno, cual corresponde a evidencia y cual corresponde a decision. Si falta alguno de estos elementos, proponga como completarlo.

\section{Presentacion convertida en practica: ANCOVA con tres variables explicativas}

Esta seccion nace de una presentacion existente y se incorpora al manual para que el material no permanezca aislado. La presentacion deja de ser un apoyo oral y se transforma en una ruta de estudio: problema, variables, datos, analisis, decision y actividad.

\begin{idea}
Ajustar una respuesta considerando covariables sin confundir control estadistico con causalidad.
\end{idea}

\subsection{Evidencia didactica extraida}

\begin{itemize}
\item Objetivo principal: comprender cómo estas características físicas explican la variabilidad en la aceptación del producto utilizando análisis de covarianza (ANCOVA).
\item La amplitud del rango en la variable de aceptación (17 puntos totales) sugiere una gran variabilidad en las respuestas de los consumidores,
\item lo cual hace especialmente valioso identificar los factores que predicen esta variación.
\end{itemize}

\subsection{Como entra al manual}

El tema se integra al capitulo relacionado con \textbf{ancova}. Si el estudiante solo observa la presentacion, ve una secuencia visual. En el manual, esa misma secuencia se expande con la pregunta de ingenieria, la razon del metodo, la interpretacion y los limites de la conclusion.

\subsection{Practica asociada}

La practica generada pide al estudiante transformar la presentacion en una decision: definir variables, organizar datos, graficar, aplicar el metodo y redactar una conclusion prudente. Si no existia practica, el sistema crea una primera version editable. Si existia una practica relacionada, esta se usa para enriquecer la narrativa y no como bloque aislado.

\subsection{Actividad de evaluacion}

Explique que parte de la presentacion corresponde a contexto, cual corresponde a diseno, cual corresponde a evidencia y cual corresponde a decision. Si falta alguno de estos elementos, proponga como completarlo.

